% Results (agent D). One section, no subsections. Every number-bearing sentence carries its source in the comment
% line above it. Real-run cells: ok rows with seeds 5660--5689 of results/tier1_summary.csv,
% results/universe_summary.csv and results/gamma_seeds/gamma_seeds_summary.csv, each run label counted once (as in
% scripts/30_talk_data.jl); W_T/W_0 = terminal / W0.

% results/tier1_summary.csv (start, stop); talk/data/market_state.csv (420 trading days)
We ran the same allocation and rebalancing system on real prices and on synthetic futures and compared the median
terminal wealth of every cell with two passive benchmarks (Tables~\ref{tab:realrun} and~\ref{tab:synthetic}). A cell
was one selection rule, one weighting rule, and one hedge setting. On real prices, the ten selection rules ($\gamma$
per sector, bandit per sector, random per sector, global $\gamma$, correlation-aware $\gamma$, liquid then $\gamma$,
random liquid 100, random whole universe, and the two hand-picked lists) and the seven weighting rules (CD, CES,
CESw, SG, MinVar, Tangent, and LS) gave $70$ cells for each hedge setting, traded over the $420$ trading days from
2 January 2025 to 4 September 2026.
% results/tier1_summary.csv, universe_summary.csv, gamma_seeds/gamma_seeds_summary.csv: 10 x 7 x 2 cells x 30 seeds
Each of the $140$ cells held $30$ runs with seeds $5660$--$5689$, for $4{,}200$ real runs, and the seed set the
controller's forecast paths and, for the random and bandit rules, the names picked.
% talk/data/largek_shares.csv; results/large_k/large_k_summary.csv (K = 100, 200)
A size comparison set global $\gamma$ against random picks from the whole universe at $K = 5$, $8$, $13$, $21$, $34$,
$100$, and $200$ names, again with $30$ runs per cell.
% results/synthetic/tables.md; results/synthetic/paths.csv (typical, burst_stress, bear_stress flags; 978 distinct
% path ids, so 22 typical paths are also stress paths)
The synthetic futures were price paths from the JumpHMM generator. We used $600$ typical paths drawn at random from
$10{,}000$ and two stress sets of $200$ paths each, the deepest SPY drawdowns in the trading window among paths with a
jump burst (burst stress) and among paths without one (bear stress), and because $22$ typical paths were also stress
paths, the three sets held $978$ distinct paths. On each path the controller recalibrated on the synthetic 2014--2024
history and traded the synthetic 2025--2026 window, and it still forecast with the multi-asset GBM while JumpHMM
supplied the prices.
% results/synthetic/rows/path_*.csv (532 run rows each: 8 deterministic rules x 7 x 2 + 6 random or bandit rules x
% 7 x 2 x 5 seeds); results/synthetic/tables.md (520,296 runs, 0 failed)
Each path carried $532$ runs, from $14$ selection rules (the ten above, plus global $\gamma$ and random whole universe
at $5$ and $34$ names), the seven weighting rules, and both hedge settings, with five seeds ($5660$--$5664$) for the
six random and bandit rules, for $520{,}296$ runs in all, none of which failed. The benchmarks were SPY buy-and-hold
and the equal-weight universe, which put $W_0/373$ into each name at the first close and never rebalanced. We
measured performance by the terminal wealth ratio $W_T/W_0$ and by continuously compounded growth rates in
$\mathrm{yr}^{-1}$. Every reported median of $W_T/W_0$ belongs to one cell and is never pooled across selection
rules, and on synthetic paths the random and bandit rules first took the median over their five seeds within each
path.

% talk/data/numbers.tex (SectorLo/Hi 1.09/1.23, BasketLo/Hi 1.03/1.70); talk/data/b3_rows_hedge_off.tex
We compared the seven weighting rules within each selection rule on real prices with the hedge off, and found that
the selection rule moved terminal wealth more than the weighting rule did (Fig.~\ref{fig:realrun};
Table~\ref{tab:realrun}). Within the $\gamma$-per-sector basket, the seven weighting rules ended between $1.086$
(MinVar) and $1.234$ (SG) times initial wealth, and the four utilities lay between $1.219$ and $1.234$, whereas
across all $70$ cells the medians ran from $1.031$ (random whole universe, with LS) to $1.697$ (hand-picked 13, with
Tangent).
% table_realrun.tex Range column and Range row (from the rounded medians); without the hand-picked rows, from the
% same table: across 0.246--0.422 (median 0.274), within 0.078--0.272 (median 0.135)
The spread within a selection rule, the highest minus the lowest of its seven medians, was $0.078$ to $0.284$, with
a median of $0.14$, and the spread across the ten selection rules within a weighting rule was $0.357$ (MinVar) to
$0.591$ (LS), with a median of $0.518$. Without the two hand-picked lists, the spread across the other eight
selection rules was $0.246$ to $0.422$ (median $0.274$), against $0.078$ to $0.272$ (median $0.135$) within them.
% talk/data/realrun_refs.csv; talk/data/realrun_grid.csv (hedge = false)
SPY buy-and-hold ended at $1.317$ and the equal-weight universe at $1.296$. Of the $70$ cells, $29$ ended above SPY,
all from five selection rules: global $\gamma$ (six of the seven weighting rules), correlation-aware $\gamma$ (five),
liquid then $\gamma$ (four), and the two hand-picked lists (seven each). None of the three per-sector rules and
neither random $13$-name rule ended above SPY with any weighting rule.
% talk/data/b3_rows_hedge_off.tex
In the $13$-name baskets, global $\gamma$ ended at $1.248$--$1.520$ against $1.031$--$1.160$ for random names from
the whole universe, whereas at three names per sector the $\gamma$ rule ($1.086$--$1.234$) and the bandit
($1.140$--$1.239$) ended close to random picks ($1.127$--$1.205$). The universe-wide rules were designed after the
per-sector results were known, so the $13$-name comparison is exploratory. The hand-picked lists ended highest, at
$1.443$--$1.697$ for the $13$ names and $1.362$--$1.646$ with WMT and PG added. Each is one list chosen by one person,
so it shows what that choice earned in this window and not the result of a repeatable rule.
% talk/data/b3_rows_hedge_off.tex (row minima)
MinVar or LS gave the lowest median under every selection rule, MinVar under four rules and LS under six.
% results/universe_summary.csv, gamma_seeds_summary.csv: quartiles of terminal / W0 over the 30 seeds, CD, hedge off
For the deterministic rules the seed changed only the controller's forecast paths, so their $30$ runs clustered
tightly (hand-picked 13 with CD: interquartile range $1.546$--$1.553$), whereas the runs of random picks spread over
a wider band (random whole universe with CD: $1.048$--$1.220$).

% talk/data/b3_rows_hedge_on.tex against b3_rows_hedge_off.tex (talk/data/realrun_grid.csv, cell medians)
We repeated every real run with the beta hedge on and compared terminal wealth, portfolio beta, and drawdown with
the hedge-off runs of the same selection and weighting rules, and found that the hedge lowered market exposure but
cost terminal wealth in most cells (Fig.~\ref{fig:realrun}; Table~\ref{tab:S-realrun-on}). At each re-plan with the
hedge on, the allocation obeyed a cap on portfolio $\beta$ that fell as the market state $\xi$ rose, and up to five
names whose online $\beta$ and down-day $\beta$ were at or below zero could enter as a hedge pool. The median
$W_T/W_0$ was lower with the hedge on in $50$ of the $70$ cells, and the median of the $70$ changes was $-0.033$.
Wealth fell in six or seven of the seven cells for the bandit rule, the three random rules, and liquid then
$\gamma$, and in five and six for the two hand-picked lists. Correlation-aware $\gamma$ ended higher in five of its
seven cells, and LS weights ended higher under six of the ten selection rules.
% results/*_summary.csv columns mean_portfolio_beta (daily sum of n_i close_i beta_i / V, averaged over days),
% mean_hedge_weight, max_drawdown (cell medians, on minus off: -0.018 to 0.019, median 0.001, above zero in 46 of 70
% cells before rounding), as in Table S-exposure; code/src/StudyBacktest.jl
We measured exposure by the portfolio $\beta$, the value-weighted online $\beta$ of the holdings divided by total
wealth, averaged over the days of a run. The hedge lowered this mean portfolio $\beta$ in all $70$ cells (cell
medians $0.388$--$0.903$ with the hedge off and $0.314$--$0.849$ with it on), and hedge-pool names held
$8.4$--$39.8\%$ of wealth on an average day (cell medians). The maximum drawdown nonetheless changed little, by at
most $0.019$ in either direction, with a median change of $+0.001$ across the $70$ cells
(Table~\ref{tab:S-exposure}).
% results/tables/hedge_effect_drawdown.csv (bandit and random rows, n_pairs = 30, win_rate <= 0.467)
When the hedge-on and hedge-off runs were paired by seed, the hedge-on run had the smaller maximum drawdown in no
more than $14$ of $30$ pairs for any weighting rule under the bandit and random per-sector rules.
% talk/data/market_state.csv (falling = 1 on 75 of 420 days); results/*_summary.csv n_bear = 75, n_bull = 344,
% spy_g_bear = 0.266, spy_g_bull = 0.144, g_bear (cell medians, on minus off); standard error of the SPY difference
% (Welch) from talk/data/market_state.csv spy column: 0.59 /yr
The cap explained part of the wealth loss. The state $\xi$ was positive on $75$ of the $420$ trading days
(Fig.~\ref{fig:market-state}), and the portfolio's mean growth rate on the $75$ days that followed a positive $\xi$
fell with the hedge on in $63$ of the $70$ cells (median change $-0.153\,\mathrm{yr}^{-1}$). On those days SPY's mean
daily growth rate was $0.266\,\mathrm{yr}^{-1}$, against $0.144\,\mathrm{yr}^{-1}$ on the other $344$ of the $419$
growth days, so a tighter cap held less of a market that rose on average. The difference of $0.12\,\mathrm{yr}^{-1}$
was small next to its standard error of $0.59\,\mathrm{yr}^{-1}$, so it describes this window and not a property of
the market state.
% results/tables/passive.csv (SPY max_drawdown 0.190)
The real window was one historical path, with a largest SPY drawdown of $0.190$, so we also tested the hedge on
simulated sell-offs.

% talk/data/mpc_events.csv, mpc_path.csv, mpc_cone_1.csv, mpc_cone_2.csv (the hand-picked 13, CD, hedge off,
% seed 5660; scripts/30_talk_data.jl MPC_LABEL); day 29 of mpc_cone_2.csv: mean 1.12346, sd 0.03553; path 1.04885
We followed the controller through one real episode and counted re-plans by trigger across the $4{,}200$
main-comparison runs, and found that the forecast band and the drawdown stop brought forward a share of re-plans
that depended on the selection rule (Fig.~\ref{fig:mpc}; Table~\ref{tab:S-replans}). The episode came from the hand-picked 13 with CD weights,
the hedge off, and seed $5660$, and ran from a re-plan on 10 October 2025 to 22 December 2025. At each re-plan the
controller forecast portfolio value over the next $21$ trading days with the multi-asset GBM and $1{,}000$ paths,
while the portfolio traded real prices. At the horizon re-plan on 10 November the portfolio stood at $1.101$ times
initial wealth, $0.61$ standard deviations (SD) above the mean of the forecast made on 10 October. Eight trading days
into the next forecast, on 20 November, the portfolio had fallen to $1.0489$ against a forecast mean of $1.1235$ with
an SD of $0.0355$, $2.10$~SD below the mean and outside the $1.96$~SD band, and the controller re-planned. The third
forecast ran to its horizon on 22 December, when the portfolio stood at $1.148$, $1.09$~SD above its mean.
% results/*_summary.csv (4,200 real runs): cell medians of (n_band_exit + n_drawdown) / (n_decisions - 1), as in
% Table S-replans (this run's cell: 0.370); per-rule ranges are hedge-off cells; n_horizon, n_band_exit, n_drawdown
% and n_basket_change sum to n_decisions - 1 in every run, and n_cash_revisit = 0 in all 4,200 runs
In the cell of this run, the median share of re-plans that came before the horizon was $0.37$. Across the $140$
cells, the median share brought forward by a band exit or a drawdown stop ranged from $0.20$ to $0.61$
(Table~\ref{tab:S-replans}). With the hedge off, the share was lowest for the three random rules, at $0.20$ to
$0.40$, and highest for the four $\gamma$ rules, whose highest cells ran from $0.53$ to $0.61$. The other re-plans
came at the horizon or on a basket change, a day on which a timer pick changed the set of names, and no run
re-planned on a cash revisit.

% results/calendar_baseline/calendar_baseline_summary.csv, each run paired with its triggered twin in
% results/*_summary.csv (4,200 pairs); cell medians over the 30 seeds of terminal / W0, max_drawdown and total_cost
% (results/calendar_baseline/calendar_vs_triggers.csv); n_horizon = 19 and every other trigger count 0 in all 4,200
% calendar-only runs; cost ratios of the cell medians 1.027--1.331 hedge off, 1.042--1.451 hedge on, as in Table
% S-replans; triggered cell medians of n_decisions - 1: 23 to 34
We repeated the main comparison with the band exit and the drawdown stop switched off, and found that the early
re-plans lowered the worst drawdown more consistently than they changed terminal wealth (Table~\ref{tab:S-calendar}).
Each of the $4{,}200$ calendar-only runs kept the selection rule, weighting rule, hedge setting and seed of one
triggered run, so the two runs of a pair differed only in the two triggers, and every calendar-only run re-planned
at the $21$-day horizon alone, $19$ times over the window, whereas in every cell the triggered runs re-planned more
often in the median. We ran this comparison because band exits measure how often the forecast missed, not whether re-planning
early helped the portfolio. With the hedge off, the median maximum drawdown was smaller with the triggers in $55$ of
the $70$ cells, by up to $0.053$ (global $\gamma$), and at most $0.015$ larger in the others. With the
hedge on it was smaller in $61$ of the $70$ cells, by up to $0.054$. Median terminal wealth was higher with the
triggers in $38$ of the $70$ cells with the hedge off and in $26$ with the hedge on, and the differences ran from
$-0.093$ to $0.191$ of initial wealth with the hedge off and from $-0.180$ to $0.183$ with it on. The sign of the
wealth difference did not follow the kind of selection rule: with the hedge off, the triggers raised terminal
wealth under all seven weighting rules for random liquid 100 and under none for random whole universe or liquid then
$\gamma$. Every triggered cell paid more trading cost than its calendar-only twin, $3$ to $33\%$ more with the hedge
off and $4$ to $45\%$ more with it on (Table~\ref{tab:S-replans}), since the triggered runs re-planned more often. On
the two stress sets of synthetic futures, the triggers likewise lowered the worst drawdown in most cells without
consistently raising terminal wealth (Tables~\ref{tab:S-calendar-burst} and~\ref{tab:S-calendar-bear}).

% talk/data/largek_shares.csv (share = fraction of the 30 x 30 run pairs with global gamma higher)
We compared the $K$ highest-$\gamma$ names in the universe with $K$ names drawn at random from the whole universe
for $K$ from $5$ to $200$ on real prices, and found that the advantage of selection by $\gamma$ shrank as the basket
grew (Fig.~\ref{fig:S-size}). For each basket size, weighting rule, and hedge setting, we formed all $900$ pairs of
one of the $30$ global $\gamma$ runs with one of the $30$ random runs, regardless of seed, and counted the share of
pairs in which global $\gamma$ ended with the higher $W_T/W_0$. Like the $13$-name comparison, this size comparison
is exploratory.
% talk/slides/b2-large-books.tex; code/src/TalkData.jl (talk_largek_shares, sg_max_K = 34);
% arxiv-paper/sections/method_allocation.tex (SG decisions infeasible at K = 100 and 200)
Stone--Geary was left out at $100$ and $200$ names, because its baseline holding of $1\%$ of wealth per name used the
whole budget at $100$ names and twice the budget at $200$, so most of its decisions were infeasible.
% talk/data/largek_shares.csv (hedge off)
With the hedge off, global $\gamma$ ended higher in $91.8$--$100\%$ of pairs at $5$ names, $94.7$--$100\%$ at $8$, and
$90.2$--$100\%$ at $13$ across the weighting rules. The share fell to $74.9$--$96.0\%$ at $21$ names and
$56.6$--$96.0\%$ at $34$, and to $33.3$--$83.2\%$ at $100$ names and $34.6$--$82.7\%$ at $200$. Random picks ended
higher in most pairs for MinVar at $100$ names and for CD, CES, and Tangent at $200$.
% talk/data/largek_shares.csv (global_median, random_median, CD, hedge off)
With CD weights, the global $\gamma$ median fell from $1.568$ at $K = 5$ to $1.520$ at $K = 13$, $1.333$ at
$K = 34$, $1.244$ at $K = 100$, and $1.185$ at $K = 200$, while the random median rose from $1.048$ to $1.098$,
$1.182$, $1.203$, and $1.211$ at the same sizes.
% talk/data/largek_shares.csv (hedge on; non-MinVar at K = 21: 94.6--100, at K = 34: 90.3--100)
With the hedge on, global $\gamma$ ended higher in $85.9$--$100\%$ of pairs at $5$ to $13$ names, and in
$90.3$--$100\%$ at $21$ and $34$ names for every weighting rule except MinVar, whose share fell to $43.2\%$ at
$21$ names and $26.1\%$ at $34$. At $100$ and $200$ names with the hedge on, the shares were $48.4$--$85.2\%$ and
$30.7$--$81.6\%$. Selection by $\gamma$ thus mattered most for baskets of $5$ to $13$ names, and at $100$ or more of
the $373$ names the CD medians of the two rules differed by less than $0.05$.

% talk/data/sf_summary.csv, sf_qq.csv, sf_acf.csv, sf_abs_acf.csv (the first 100 typical paths by path id,
% scripts/30_talk_data.jl N_SYN_PATHS; market factor g_m over 2014--2024 against real SPY over 2014--2024; kurtosis
% is non-excess, 3 for a normal; abs ACF at lag 21: real 0.134, generator median 0.0067)
We compared daily growth rates of the JumpHMM market generator with those of real SPY over 2014--2024 and checked
the composed multi-asset paths against the real universe, and found that the generator reproduced fat tails but not
the persistence of volatility (Fig.~\ref{fig:jumphmm}; Table~\ref{tab:S-checks}). For the comparison with SPY we
used the market factor of the first $100$ typical paths by path number. Across these paths the median kurtosis of
the daily growth rate was $13.7$, against $16.8$ for SPY and $3$ for a normal distribution. In the tails the median
generator quantile lay at or beyond SPY's, at $-17.2$ against $-16.3\,\mathrm{yr}^{-1}$ for the $0.1\%$ quantile,
$-10.1$ against $-7.8$ for the $1\%$ quantile, $7.3$ against $6.4$ for the $99\%$ quantile, and $15.1$ against
$14.8$ for the $99.9\%$ quantile. Both series had small autocorrelation of $g$ at lag 1 ($-0.04$ for the generator
median, $-0.11$ for SPY). The autocorrelation of $|g|$ differed, since for the generator median it fell from $0.22$
at lag 1 to $0.08$ at lag 2, $0.03$ at lag 3, and $0.007$ at lag 21, whereas SPY's was $0.36$ at lag 1 and still
$0.13$ at lag 21, about one trading month.
% results/synthetic/checks.md (first 50 typical paths, recalibrated on synthetic 2014--2024; V1 over all 10,000)
Of the $22$ gated checks on the composed paths, each with a pass tolerance, $21$ passed. The median synthetic SPY
growth rate over 2014--2024 was $0.105\,\mathrm{yr}^{-1}$ across the $10{,}000$ paths, against
$0.106\,\mathrm{yr}^{-1}$ for real SPY. The median name's variance was $0.994$ of the variance of its own JumpHMM
draw, and $1.066$ of its real variance, a comparison reported without a tolerance. The mean residual correlation
within each of the 11 sectors lay within $0.05$ of the real value (for example $0.535$ against $0.542$ in Energy),
and the SD of the online $\beta$ was $0.985$ of the real one. The failing check was the cross-name SD of $\alpha$,
the intercept of each name's least-squares regression on the market over 2014--2024, which was
$0.089\,\mathrm{yr}^{-1}$ against $0.076\,\mathrm{yr}^{-1}$ for the real universe, outside the tolerance of
$0.01\,\mathrm{yr}^{-1}$.
% results/synthetic/checks.md (V5 across sectors; V4 synthetic - real at lags 1 and 5, synthetic - own draw at lag 1)
Comparisons without a tolerance showed what the composition left out. The mean residual correlation across sectors
was $0.0002$ against $0.031$, because cross-sectional dependence entered only through the market and sector factors.
The median name's autocorrelation of $|g|$ was $0.10$ below the real one at lag 1 and $0.17$ below at lag 5, but
within $0.04$ of its own JumpHMM draw, so the shortfall came from the per-name draws and not from the composition.
% results/synthetic/tables.md lines 9 and 1067 (names whose online beta went to 0 or below: synthetic 2025--26
% median 7 per path, range 1 to 19; real 2014--2024 windows of the same length, median 5, range 0 to 53 over 112
% windows; real 2025--26: 120)
The synthetic paths also offered the hedge fewer candidates, with a median of seven names per path whose online
$\beta$ reached zero or below in the 2025--2026 window, against $120$ in the real 2025--2026 window and a median of
five in real 2014--2024 windows of the same length (range $0$ to $53$). The synthetic paths thus resembled a typical
stretch of 2014--2024, and the real trading window held more hedge candidates than any earlier stretch of that
length.

% talk/data/syn_typical_refs.csv, syn_typical_dots.csv; results/synthetic/tables/typical_hedge_off_wealth_median.csv
We ran all selection and weighting rules on the $600$ typical futures, compared selection by $\gamma$ with random
picks of the same size, and placed each real result among the typical futures of its cell, and found that selection
by $\gamma$ beat random picks in most typical futures (Tables~\ref{tab:synthetic} and~\ref{tab:S-head-to-head};
Fig.~\ref{fig:S-percentile}). In the typical futures the median synthetic SPY ended at $1.185$ and the equal-weight
universe at $1.230$, and the $70$ hedge-off cell medians of the ten main selection rules ran from $1.097$ (random
whole universe, with LS) to $1.244$ (liquid then $\gamma$, with SG).
% results/synthetic/tables/typical_hedge_off_wealth_median.csv and typical_hedge_off_wealth_beat_spy.csv (the ten
% selection rules of Fig. fig:synthetic)
Of these $70$ cells, $45$ ended above the SPY median and seven above the equal-weight median. The equal-weight
universe ended above its own path's SPY in $77.5\%$ of typical futures, whereas on real prices it ended below SPY
($1.296$ against $1.317$).
% results/synthetic/tables/typical_head_to_head.csv (wealth, share_above_0 and median)
With the hedge off, global $\gamma$ ended above $13$ random names from the whole universe in $62.3$--$71.2\%$ of
futures across the weighting rules, with a median difference in $W_T/W_0$ of $0.042$ to $0.101$, and $\gamma$ per
sector ended above random per sector in $60.0$--$73.7\%$. With the hedge on, the $13$-name share was
$63.2$--$71.5\%$.
% talk/data/syn_typical_dots.csv (hedge off; the four gamma rules: gamma per sector, global gamma, correlation-aware
% gamma, liquid then gamma)
The hand-picked 13 ended between $1.175$ and $1.242$, within the $1.142$--$1.244$ range of the four $\gamma$ rules
($\gamma$ per sector, global $\gamma$, correlation-aware $\gamma$, and liquid then $\gamma$), whereas on real prices
both hand-picked lists ended highest.
% talk/slides/b1-percentiles.tex (footer); code/src/SyntheticTables.jl (real_percentiles, path_cells);
% results/synthetic/tables/typical_real_percentile_hedge_off.csv
To place the real results, we used the seeds of the synthetic runs, $5660$ for the deterministic rules and the median
over $5660$--$5664$ for the random and bandit rules, rather than the $30$-seed medians, and took the percentile of
each real value among the $600$ typical futures of its cell. Real SPY sat at percentile $70$ of synthetic SPY and the
real equal-weight universe at $60.7$, so the real window was a better-than-median period for the market, and part of
any real margin over the typical futures reflected that period. The real results of the hand-picked 13 sat at
percentiles $79.2$ to $90.7$ of their cells, those of global $\gamma$ at $47.5$ to $81.7$, those of $\gamma$ per
sector at $39.8$ to $56.2$, and those of $13$ random names from the whole universe at $43.3$ to $64.3$. Of the $70$
hedge-off cells, $52$ sat above the median future and $19$ above SPY's percentile of $70$.
% results/synthetic/tables/typical_real_percentile_hedge_off.csv (hand-picked rows: 14 of 14 above 70, lowest 74.0;
% the four gamma rules: 5 of 28 above 70)
The hand-picked lists sat at or above percentile $74$ in all $14$ of their hedge-off cells, above the market's $70$,
whereas the four $\gamma$ rules sat at percentiles $39.8$ to $81.7$ and above $70$ in five of their $28$ cells
(Fig.~\ref{fig:S-percentile}). In the real window the hand-picked lists thus exceeded their typical futures by more
than the market did, and the $\gamma$ rules in most cells did not.

% results/synthetic/paths.csv (window_max_drawdown and window_growth, medians per set)
We repeated the comparison on the two stress sets of simulated sell-offs, and found that every selection and
weighting median ended above SPY and the equal-weight universe (Fig.~\ref{fig:synthetic};
Table~\ref{tab:synthetic}). The median synthetic SPY drawdown in the 2025--2026 window was $0.364$ in the burst
stress set and $0.406$ in the bear stress set, against $0.176$ in the typical futures, and SPY's median growth rate
was $-0.087$ and $-0.199\,\mathrm{yr}^{-1}$, against $0.102\,\mathrm{yr}^{-1}$.
% talk/data/syn_burst_stress_dots.csv, _refs.csv; talk/data/syn_bear_stress_dots.csv, _refs.csv;
% results/synthetic/tables/{burst,bear}_stress_hedge_{off,on}_wealth_median.csv (all 14 rules, 98 cells per setting)
In both stress sets all $140$ cell medians of the ten main selection rules ended above the SPY median and the
equal-weight median, and so did all $98$ cells per hedge setting of the $14$ synthetic selection rules. Among the
ten main selection rules, the lowest median in the burst stress set was $0.924$ (random per sector, MinVar, hedge off), against $0.865$ for SPY and $0.916$
for the equal-weight universe, and the highest was $1.031$ (hand-picked 13, LS, hedge off). In the bear stress set
the lowest median was $0.774$ (random whole universe, Tangent, hedge off), against $0.718$ for SPY and $0.753$ for
the equal-weight universe, and the highest was $1.001$ (hand-picked 13 with WMT and PG, LS, with either hedge
setting).
% results/synthetic/tables/{burst,bear}_stress_hedge_{off,on}_wealth_beat_spy.csv
Each cell also ended above its own path's SPY in $69.0$--$88.5\%$ of burst futures and $72.5$--$100\%$ of bear
futures. Random picks also ended above SPY, so the margin over SPY did not come from selection by $\gamma$.
% results/synthetic/tables/{burst,bear}_stress_hedge_off_wealth_median.csv (SPY 0.86496 and 0.71847, equal weight
% 0.91555 and 0.75313, lowest cells 0.92412 and 0.77386)
The equal-weight universe, which neither caps volatility nor re-plans, also ended above SPY, by $0.051$ in the burst
stress set and $0.035$ in the bear stress set, and the lowest cell medians exceeded it by $0.009$ and $0.021$, so part
of the margin over SPY came from the synthetic market and the rest from parts of the system that every cell shared.
% results/synthetic/rows/path_*.csv mean_portfolio_beta, ten main selection rules, hedge off: median over seeds within
% a path, then over the 200 paths of each stress set, as in Table S-synbeta
With the hedge off, the cell medians of the mean portfolio $\beta$ were $0.322$ to $0.764$ in the burst stress set and
$0.229$ to $0.850$ in the bear stress set, all below one (Table~\ref{tab:S-synbeta}).
% results/synthetic/tables/bear_stress_head_to_head.csv, burst_stress_head_to_head.csv (wealth, share_above_0)
In bear futures selection by $\gamma$ lost its typical-future advantage for five of the seven weighting rules. With
the hedge off, global $\gamma$ ended above $13$ random names in $28.0$--$38.0\%$ of bear futures for CD, CES, CESw,
SG, and LS, and at $5$ names the same five weighting rules gave $38.5$--$45.0\%$. MinVar and Tangent, whose weights do
not use $\gamma$, reversed this pattern, with global $\gamma$ ahead in $70.0\%$ and $71.0\%$ of bear futures. In
burst futures global $\gamma$ ended above $13$ random names in $52.5$--$67.5\%$ of futures with the hedge off.
% talk/data/syn_{typical,burst_stress,bear_stress}_dots.csv (hedge on minus hedge off, per cell)
With the hedge on, the median $W_T/W_0$ ended higher in $65$ of the $70$ cells of the bear stress set (median of the
$70$ cell changes $+0.029$), and the lowest cell median rose from $0.774$ to $0.814$. In the burst stress set it
ended higher in $34$ of the $70$ cells (median change $-0.001$), and in the typical futures in $11$ of the $70$
(median change $-0.006$). The $\beta$ cap and hedge pool thus raised terminal wealth in prolonged declines and
lowered it in typical futures, as they did on real prices.

% results/calendar_baseline/synthetic_vs_triggers.csv (scripts/dev/calendar_baseline_synthetic.jl): within a path the
% median over seeds of each cell, then the median across the 200 paths of the set. Triggered minus calendar-only cell
% medians: max_drawdown -0.0222 to 0.0073 (median of the 280 cell reductions 0.0091; real prices 0.0111 over 140
% cells, results/calendar_baseline/calendar_vs_triggers.csv), terminal / W0 -0.0309 to 0.0287. Calendar-only: 19
% re-plans in every run; lowest cell medians 0.93009 (burst) and 0.77246 (bear). Triggered cell medians of
% n_decisions - 1: 25 to 32
We repeated the calendar-only comparison on the $400$ stress futures, and found that the early re-plans again lowered
the worst drawdown more consistently than they changed terminal wealth, by about $0.01$ in the median cell as on real
prices, although the stress declines were deeper (Tables~\ref{tab:S-calendar-burst} and~\ref{tab:S-calendar-bear}).
The median maximum drawdown was smaller with the triggers in $65$ of the $70$ cells with the hedge off and in all $70$
with it on in the burst stress set, and in $65$ and $67$ of the $70$ in the bear stress set, by up to $0.022$, and at
most $0.007$ larger in the others. Median terminal wealth was higher with the triggers in $33$ and $28$ of the $70$
cells in the burst stress set and in $51$ and $56$ in the bear stress set, with differences from $-0.031$ to $0.029$
of initial wealth. Every calendar-only run re-planned $19$ times, every $21$ trading days, while the triggered cells
re-planned $25$ to $32$ times in the median. With the two triggers switched off, every cell median still ended above
the SPY median and the equal-weight median in both stress sets, the lowest at $0.930$ in the burst stress set and
$0.772$ in the bear stress set, so the margin over SPY did not depend on the early re-plans.

